% =====================================================================
\chapter{Regresión lineal y mínimos cuadrados}
\label{tema:2}
% =====================================================================

Este tema estudia el modelo más sencillo de aprendizaje supervisado, la regresión lineal ajustada por mínimos cuadrados. Se deducen las ecuaciones normales, se interpretan como una proyección ortogonal y se justifica la pérdida cuadrática como máxima verosimilitud con ruido gaussiano. La exposición sigue el capítulo~1 de los apuntes de CS229~\cite{cs229}, las secciones~4.2 y~4.3 de Strang~\cite{strang} y el curso 18.06SC del MIT~\cite{mit1806}, la sección~3.1 de Bishop~\cite{bishop} y la sección~3.1 de Zhang \emph{et al.}~\cite{d2l}.

\section{El modelo lineal y la matriz de diseño}

\begin{definicion}[Regresión lineal]
Un \emph{modelo de regresión lineal} predice $\hat y=w^{\top}x+b$, con \emph{pesos} $w\in\R^{d}$ y \emph{término independiente} $b\in\R$. Añadiendo a cada entrada una componente constante, $\tilde x=(1,x)\in\R^{d+1}$, y escribiendo $\theta=(b,w)$, el modelo queda $\hat y=\theta^{\top}\tilde x$. En lo sucesivo se omite la tilde.
\end{definicion}

\begin{definicion}[Matriz de diseño y suma de cuadrados]
Dados $n$ ejemplos, la \emph{matriz de diseño} $X\in\R^{n\times(d+1)}$ tiene por fila $i$ el vector $x_i^{\top}$ (con su primera componente igual a~1), y $y=(y_1,\dots,y_n)^{\top}$. El vector de predicciones es $\hat y=X\theta$, y la \emph{suma de cuadrados de los residuos} es
\[
J(\theta)=\norm{y-X\theta}^{2}=\sum_{i=1}^{n}\bigl(y_i-\theta^{\top}x_i\bigr)^{2}.
\]
El estimador de \emph{mínimos cuadrados} es cualquier $\hat\theta$ que minimice $J$; se escribe $\SSE=J(\hat\theta)$.
\end{definicion}

\begin{observacion}[Funciones base]
El modelo es lineal en $\theta$, no necesariamente en $x$. Si se sustituye $x$ por un vector de \emph{funciones base} $\phi(x)$, por ejemplo $\phi(x)=(1,x,x^{2},\dots,x^{p})$, todo lo que sigue se aplica sin cambios con $X_{ij}=\phi_j(x_i)$~\cite[§3.1]{bishop}. Así se ajustan los polinomios del \cref{tema:5}.
\end{observacion}

\section{Las ecuaciones normales}

\begin{lema}[Gradientes de formas lineales y cuadráticas]\label{lem:gradientes}
Sean $a\in\R^{m}$ y $A\in\R^{m\times m}$. Entonces $\nabla_\theta(a^{\top}\theta)=a$ y $\nabla_\theta(\theta^{\top}A\theta)=(A+A^{\top})\theta$, que vale $2A\theta$ si $A$ es simétrica.
\end{lema}

\begin{teorema}[Ecuaciones normales]\label{teo:normales}
Un vector $\hat\theta$ minimiza $J(\theta)=\norm{y-X\theta}^{2}$ si y solo si cumple las \emph{ecuaciones normales}
\[
X^{\top}X\,\hat\theta=X^{\top}y.
\]
Además, para cualquier solución $\hat\theta$ y cualquier $\theta$, se verifica $J(\theta)=J(\hat\theta)+\norm{X(\theta-\hat\theta)}^{2}$.
\end{teorema}

\begin{lema}\label{lem:rango}
$X^{\top}X$ es simétrica y semidefinida positiva. Es definida positiva, y por tanto invertible, si y solo si las columnas de $X$ son linealmente independientes, es decir, si $\rg(X)=d+1$.
\end{lema}

\begin{corolario}[Unicidad]\label{cor:unicidad}
Si $\rg(X)=d+1$, el estimador de mínimos cuadrados es único y vale $\hat\theta=(X^{\top}X)^{-1}X^{\top}y$. Si $\rg(X)<d+1$, hay infinitos minimizadores, todos con las mismas predicciones $X\hat\theta$, y el de norma mínima es $X^{+}y$, donde $X^{+}$ es la pseudoinversa (módulo~3).
\end{corolario}

\addtocounter{definicion}{1}%

\begin{observacion}[Cálculo numérico]
En la práctica no se calcula $(X^{\top}X)^{-1}$: se resuelve el sistema con una factorización de Cholesky, o se aplica a $X$ una factorización QR o la descomposición en valores singulares, que son más estables porque el número de condición de $X^{\top}X$ es el cuadrado del de $X$. Con millones de parámetros, como en una red neuronal, no existe fórmula cerrada y se usa el descenso de gradiente (\cref{sec:lms}).
\end{observacion}

\section{Interpretación geométrica}

Las ecuaciones normales equivalen a $X^{\top}(y-X\hat\theta)=0$: el residuo es ortogonal a todas las columnas de $X$. Por tanto, $\hat y=X\hat\theta$ es el punto del espacio de columnas $\mathcal{C}(X)$ más próximo a $y$, su \emph{proyección ortogonal}~\cite[§4.2--4.3]{strang}.

\begin{proposicion}[Matriz de proyección]\label{prop:proyeccion}
Supóngase que $\rg(X)=d+1$ y sea $P=X(X^{\top}X)^{-1}X^{\top}$, de tamaño $n\times n$. Entonces:
\begin{enumerate}
  \item $\hat y=Py$ y $r=(I-P)y$;
  \item $P^{\top}=P$ y $P^{2}=P$;
  \item $\hat y^{\top}r=0$ y, en consecuencia, $\norm{y}^{2}=\norm{\hat y}^{2}+\norm{r}^{2}$;
  \item si $X$ contiene la columna de unos, $\sum_i r_i=0$, y la recta (o el hiperplano) ajustada pasa por el punto medio $(\bar x,\bar y)$.
\end{enumerate}
\end{proposicion}

\addtocounter{figure}{1}%

En el \cref{ej:recta}, $\norm{y}^{2}=46$, $\norm{\hat y}^{2}=45{,}8$ y $\norm{r}^{2}=0{,}2$, de acuerdo con la \cref{prop:proyeccion}(iii).

\section{Descenso de gradiente: el algoritmo LMS}
\label{sec:lms}

Cuando $n$ o $d$ son grandes, se minimiza $J$ iterativamente. Con la pérdida media $\frac{1}{2n}\norm{y-X\theta}^{2}$, el gradiente es $-\frac1n X^{\top}(y-X\theta)$, y el descenso de gradiente con tasa $\alpha>0$ actualiza
\[
\theta\leftarrow\theta+\frac{\alpha}{n}\sum_{i=1}^{n}\bigl(y_i-\theta^{\top}x_i\bigr)x_i.
\]
Si en cada paso se usa un solo ejemplo, se obtiene la \emph{regla LMS} (\emph{least mean squares}) o de Widrow-Hoff, $\theta\leftarrow\theta+\alpha\,(y_i-\theta^{\top}x_i)\,x_i$: la corrección es proporcional al error cometido en ese ejemplo~\cite[§1.1]{cs229}. Como $J$ es convexa (su hessiano $2X^{\top}X$ es semidefinido positivo, \cref{lem:rango}), el descenso de gradiente converge al mínimo global si la tasa es suficientemente pequeña (módulo~7).

\section{Interpretación probabilística}

La pérdida cuadrática no es una elección arbitraria: es la que resulta de suponer ruido gaussiano~\cite[§1.3]{cs229}, \cite[§3.1.1]{bishop}.

\begin{definicion}[Modelo lineal gaussiano]
Las respuestas cumplen $y_i=\theta^{\top}x_i+\varepsilon_i$, donde los errores $\varepsilon_1,\dots,\varepsilon_n$ son independientes con distribución $\N(0,\sigma^{2})$. Equivalentemente, $y_i\mid x_i\sim\N(\theta^{\top}x_i,\sigma^{2})$.
\end{definicion}

\begin{teorema}[Mínimos cuadrados y máxima verosimilitud]\label{teo:mv-gauss}
En el modelo lineal gaussiano, la log-verosimilitud negativa es
\[
-\ln L(\theta,\sigma^{2})=\frac{n}{2}\ln(2\pi\sigma^{2})+\frac{1}{2\sigma^{2}}\norm{y-X\theta}^{2}.
\]
Para cualquier $\sigma^{2}>0$, el estimador de máxima verosimilitud de $\theta$ es el de mínimos cuadrados, y el de $\sigma^{2}$ es $\hat\sigma^{2}=\SSE/n$.
\end{teorema}

En el \cref{ej:recta}, $\hat\sigma^{2}=0{,}2/4=0{,}05$. El razonamiento se generaliza: cada supuesto sobre la distribución de $y$ produce su pérdida de máxima verosimilitud (\cref{tab:mv}). Con ruido de Laplace, por ejemplo, la pérdida es la suma de valores absolutos (\cref{eje:2-laplace}).

\begin{table}[h]
\centering\small
\caption{Pérdida de máxima verosimilitud según la distribución supuesta para $y$.}\label{tab:mv}
\begin{tabular}{@{}llll@{}}
\toprule
Respuesta & Distribución & Pérdida & Modelo \\
\midrule
real & normal & cuadrática & regresión lineal (\cref{tema:2}) \\
binaria & Bernoulli & entropía cruzada binaria & regresión logística (\cref{tema:3}) \\
una de $K$ clases & categórica & entropía cruzada & regresión softmax (\cref{tema:4}) \\
token siguiente & categórica sobre $\V$ & entropía cruzada & capa de salida de un LLM \\
\bottomrule
\end{tabular}
\end{table}

\begin{observacion}[Física]
$\frac12\sum_i(y_i-\hat y_i)^{2}$ es la energía elástica de $n$ muelles verticales de constante unidad que unen cada dato con la recta. La recta de mínimos cuadrados es la posición de equilibrio, y las ecuaciones normales expresan que la fuerza total y el momento total de los muelles son nulos.
\end{observacion}

